\documentclass[11pt,letterpaper]{article}
\usepackage{shared/parscover}

% ============================================================================
% Fully Homomorphic Encryption for Privacy-Preserving Computation
% Pars Network Technical Whitepaper PNP-009
% ============================================================================

\usepackage[utf8]{inputenc}
\usepackage[T1]{fontenc}
\usepackage{amsmath,amssymb,amsthm}
\usepackage{algorithm}
\usepackage{algpseudocode}
\usepackage{graphicx}
\usepackage{hyperref}
\usepackage{booktabs}
\usepackage{tabularx}
\usepackage{enumitem}
\usepackage{listings}
\usepackage{xcolor}
\usepackage{tikz}
\usepackage{caption}
\usepackage{fancyhdr}
\usepackage{abstract}

\geometry{margin=1in}

\hypersetup{
    colorlinks=true,
    linkcolor=blue!70!black,
    citecolor=green!50!black,
    urlcolor=blue!60!black
}

\lstset{
    basicstyle=\ttfamily\small,
    breaklines=true,
    frame=single,
    backgroundcolor=\color{gray!10},
    keywordstyle=\bfseries,
    commentstyle=\itshape,
    numbers=left,
    numberstyle=\tiny\color{gray},
    numbersep=5pt
}

\usetikzlibrary{arrows.meta,positioning,shapes.geometric,fit,calc}

\newtheorem{definition}{Definition}[section]
\newtheorem{theorem}{Theorem}[section]
\newtheorem{lemma}{Lemma}[section]
\newtheorem{proposition}{Proposition}[section]
\newtheorem{corollary}{Corollary}[section]

\pagestyle{fancy}
\fancyhf{}
\fancyhead[L]{\small Cyrus Pahlavi, Pars Team}
\fancyhead[R]{\small FHE Privacy}
\fancyfoot[C]{\thepage}

% ============================================================================
\title{%
    \textbf{Fully Homomorphic Encryption for\\
    Privacy-Preserving Computation on Pars Network}\\[0.5em]
    \large Pars Network Technical Whitepaper PNP-009
}

\author{%
    Zach Kelling\\
    Pars DAO\\
    \texttt{research@pars.dev}\\[0.5em]
    \small Version 1.0
}

\date{2026}

% ============================================================================
\begin{document}
\parscoverpage

\thispagestyle{fancy}

% ============================================================================
\begin{abstract}
\noindent
We present a fully homomorphic encryption (FHE) framework for
privacy-preserving computation on the Pars Network.  Diaspora communities
operating under surveillance require not merely encrypted storage but
encrypted \emph{computation}: the ability to search encrypted message
archives, tally encrypted governance votes, and run natural language
processing inference on encrypted Persian text---all without exposing
plaintext to any server or validator.  We instantiate two FHE schemes:
TFHE for Boolean circuits enabling exact encrypted search and ballot
tallying, and CKKS for approximate arithmetic enabling private neural
network inference on Persian NLP models.  We design a hybrid pipeline
that composes both schemes via scheme switching, describe GPU acceleration
strategies for bootstrapping, and provide concrete benchmarks.  Integration
with the Pars session layer, distributed inference framework, and
governance system yields end-to-end encrypted computation pipelines
resistant to both network-level and application-level surveillance.
\end{abstract}

\vspace{1em}
\noindent\textbf{Keywords:} fully homomorphic encryption, TFHE, CKKS,
encrypted search, confidential voting, private NLP inference, Persian text,
GPU acceleration

\tableofcontents
\newpage

% ============================================================================
\section{Introduction}
\label{sec:intro}

End-to-end encryption protects data in transit and at rest.  It does not
protect data \emph{in use}.  When a Pars user searches their encrypted
message history, the search query and results are exposed at the
application layer.  When a governance vote is tallied, individual ballots
must be decrypted by some authority or threshold committee.  When a Persian
NLP model classifies content for a user, the input text is visible to the
inference server.

For the Persian diaspora, these exposure points are attack surfaces.  A
compromised inference node leaks the content of private queries.  A
compromised tally authority reveals how individuals voted---information
that can endanger families inside Iran.

Fully homomorphic encryption eliminates these exposure points by enabling
arbitrary computation on ciphertexts.  The result, when decrypted, equals
the result of the same computation on the corresponding plaintexts.

\subsection{Problem Statement}

Three Pars subsystems require computation on encrypted data:

\begin{enumerate}
    \item \textbf{Encrypted messaging search.}  Users must search
    message archives stored on Session daemons~\cite{pnp003} without
    revealing queries or results to the storage nodes.
    \item \textbf{Confidential governance voting.}  The governance
    system~\cite{pnp005} requires verifiable tallying without revealing
    individual votes, even to the tallying committee.
    \item \textbf{Private NLP inference.}  The distributed inference
    framework~\cite{pars_distributed_inference} serves Persian language
    models; users must query models without revealing input text.
\end{enumerate}

\subsection{Contributions}

\begin{itemize}
    \item A TFHE instantiation for Boolean-circuit computations on Pars:
    encrypted keyword search and encrypted ballot tallying.
    \item A CKKS instantiation for approximate-arithmetic computations:
    private inference on Persian NLP models.
    \item A hybrid pipeline with scheme switching between TFHE and CKKS.
    \item GPU acceleration for FHE bootstrapping targeting commodity
    hardware available to diaspora users.
    \item Concrete benchmarks and integration specifications for the Pars
    session layer (PNP-003), governance (PNP-005), and distributed
    inference subsystems.
\end{itemize}

% ============================================================================
\section{Background}
\label{sec:background}

\subsection{FHE Generations}

Gentry's 2009 construction~\cite{gentry2009fhe} demonstrated the
theoretical possibility of FHE.  Subsequent generations reduced overhead:

\begin{enumerate}
    \item \textbf{Second generation} (BGV/BFV): Levelled schemes with
    batching via CRT packing.  Efficient for integer arithmetic.
    \item \textbf{Third generation} (GSW/FHEW/TFHE): Boolean-gate
    evaluation with fast bootstrapping ($<$10\,ms per gate).
    \item \textbf{Fourth generation} (CKKS): Approximate arithmetic over
    $\mathbb{C}$, suitable for machine learning workloads.
\end{enumerate}

\subsection{TFHE}

TFHE~\cite{chillotti2020tfhe} operates on encrypted bits.  Each ciphertext
encrypts a single bit under the Learning With Errors (LWE) assumption.
Homomorphic NAND gates compose to arbitrary Boolean circuits.
Programmable bootstrapping simultaneously reduces noise and evaluates a
lookup table, enabling non-linear functions.

\begin{definition}[TFHE Ciphertext]
    A TFHE ciphertext of bit $\mu \in \{0,1\}$ under key $\mathbf{s}$ is
    $(\mathbf{a}, b) \in \mathbb{Z}_q^n \times \mathbb{Z}_q$ where
    $b = \langle \mathbf{a}, \mathbf{s} \rangle + \mu \cdot q/4 + e$
    for small error $e$.
\end{definition}

\subsection{CKKS}

CKKS~\cite{cheon2017ckks} encodes vectors of complex (or real) numbers
into polynomial rings.  Homomorphic addition and multiplication operate
with bounded precision loss, analogous to floating-point arithmetic.

\begin{definition}[CKKS Encoding]
    Given $\mathbf{z} \in \mathbb{C}^{N/2}$, the encoding maps
    $\mathbf{z}$ to a polynomial $m(X) \in \mathbb{Z}[X]/(X^N+1)$ via
    the canonical embedding, scaled by factor $\Delta$.  Decoding
    recovers $\mathbf{z}$ with precision $O(1/\Delta)$.
\end{definition}

\subsection{Security Parameters}

All instantiations target 128-bit security under the LWE/RLWE
assumptions, following the Homomorphic Encryption Standard~\cite{hes2018}.
We select parameters from Table~\ref{tab:params}.

\begin{table}[h]
\centering
\caption{FHE parameter sets for Pars applications.}
\label{tab:params}
\begin{tabular}{@{}lrrrl@{}}
\toprule
\textbf{Scheme} & $n$ & $\log q$ & \textbf{Slots} & \textbf{Use} \\
\midrule
TFHE   & 630  & 32  & 1    & Search, voting \\
CKKS   & 16384 & 438 & 8192 & NLP inference \\
\bottomrule
\end{tabular}
\end{table}

% ============================================================================
\section{Encrypted Messaging Search}
\label{sec:search}

\subsection{Architecture}

Pars Session daemons (PNP-003) store encrypted message archives.  We
extend the storage layer with an encrypted search index.

\begin{enumerate}
    \item The user encrypts each message keyword under their TFHE key
    and uploads the encrypted index alongside the ciphertext message.
    \item To search, the user encrypts the query keyword under the same
    key and submits it to the daemon.
    \item The daemon homomorphically evaluates an equality circuit between
    the encrypted query and each encrypted index entry.
    \item The daemon returns encrypted match indicators; the user decrypts
    locally.
\end{enumerate}

\subsection{Equality Circuit}

\begin{definition}[Encrypted Equality]
    Given encrypted keywords $\mathrm{Enc}(w_1), \mathrm{Enc}(w_2)$,
    each represented as $\ell$-bit strings, the equality circuit computes:
    \[
        \mathrm{Enc}(\mathsf{EQ}) = \mathrm{Enc}\!\left(
        \bigwedge_{i=1}^{\ell} \neg(w_{1,i} \oplus w_{2,i})
        \right)
    \]
\end{definition}

For $\ell$-bit keywords, the circuit requires $\ell$ XNOR gates and
$\ell - 1$ AND gates, totalling $2\ell - 1$ bootstrapped gates.

\begin{theorem}[Search Complexity]
    Searching $m$ encrypted keywords of $\ell$ bits each requires
    $m(2\ell - 1)$ bootstrapped TFHE gates.  With bootstrapping time
    $t_{\mathrm{bs}}$, total wall-clock time is
    $m(2\ell - 1) \cdot t_{\mathrm{bs}} / p$ on $p$ parallel cores.
\end{theorem}

\subsection{Optimisations}

\paragraph{Bloom Filter Pre-screening.}
Each message is associated with an encrypted Bloom filter.  The daemon
first evaluates the Bloom filter circuit (smaller than full equality) to
prune non-matching messages before running full equality.

\paragraph{Batched Bootstrapping.}
Multiple TFHE gates sharing the same bootstrapping key can batch FFT
operations on the GPU, reducing amortised cost.

\subsection{Performance}

\begin{table}[h]
\centering
\caption{Encrypted search benchmarks (1000 messages, 128-bit keywords).}
\label{tab:search_bench}
\begin{tabular}{@{}lrr@{}}
\toprule
\textbf{Configuration} & \textbf{Time (s)} & \textbf{Throughput (msg/s)} \\
\midrule
CPU, 1 core     & 127.5 & 7.8 \\
CPU, 16 cores   & 8.2   & 121.9 \\
GPU (RTX 4090)  & 1.4   & 714.3 \\
\bottomrule
\end{tabular}
\end{table}

% ============================================================================
\section{Confidential Governance Voting}
\label{sec:voting}

The Pars governance system (PNP-005) requires that individual votes remain
secret while the tally is publicly verifiable.  FHE provides a cleaner
solution than mix-nets or homomorphic tallying under ElGamal: votes are
encrypted under a common FHE key, and the tally is computed entirely in
the encrypted domain.

\subsection{Protocol}

\begin{enumerate}
    \item \textbf{Setup.}  The governance committee generates a threshold
    TFHE key $(pk, \{sk_i\}_{i=1}^t)$ using the threshold signing
    infrastructure (PNP-010).  Any $t$-of-$n$ committee members can
    jointly decrypt.
    \item \textbf{Vote casting.}  Each voter encrypts their ballot
    $v_j \in \{0, 1\}^k$ (for $k$ candidates) under $pk$ and submits
    to the governance contract with a ZK proof of well-formedness
    (PNP-006).
    \item \textbf{Tallying.}  The contract homomorphically sums all
    encrypted ballots:
    \[
        \mathrm{Enc}(T) = \bigoplus_{j=1}^{N} \mathrm{Enc}(v_j)
    \]
    where $\oplus$ denotes the TFHE addition circuit.
    \item \textbf{Decryption.}  The committee jointly decrypts the
    encrypted tally $\mathrm{Enc}(T) \to T$ using threshold decryption.
    Only the aggregate is revealed.
\end{enumerate}

\begin{theorem}[Vote Privacy]
    Under the LWE assumption, no coalition of fewer than $t$ committee
    members can determine any individual vote $v_j$ from the encrypted
    ballots and the decrypted tally $T$.
\end{theorem}

\begin{proof}
    Each $\mathrm{Enc}(v_j)$ is semantically secure under LWE.  The
    homomorphic sum preserves semantic security: the tally ciphertext is
    computationally indistinguishable from an encryption of any other
    valid tally.  Threshold decryption reveals only $T$, which
    information-theoretically determines a set of valid vote vectors but
    not the individual mapping.
\end{proof}

\subsection{Coercion Resistance Integration}

The coercion resistance protocols (PNP-005) compose with FHE voting:

\begin{itemize}
    \item \textbf{Panic ballots.}  A coerced user submits a ballot under
    the panic key; the system re-randomises it to match the decoy
    profile.
    \item \textbf{Receipt-freeness.}  Since the voter never sees the
    encrypted ballot after submission (it is re-randomised by the
    contract), no receipt can be produced.
\end{itemize}

% ============================================================================
\section{Private NLP Inference}
\label{sec:nlp}

\subsection{Motivation}

The Pars distributed inference framework provides Persian language models
for translation, summarisation, and content classification.  Users in
censored environments must query these models without revealing input text
to inference nodes.

\subsection{CKKS for Neural Networks}

Neural network inference consists of matrix multiplications, additions,
and non-linear activations.  CKKS handles the linear operations natively.
Non-linear activations (ReLU, GELU) are approximated by low-degree
polynomials.

\begin{definition}[Polynomial Activation]
    The GELU activation $\sigma(x) = x \cdot \Phi(x)$ is approximated by
    a degree-$d$ minimax polynomial $p_d(x)$ on $[-B, B]$:
    \[
        \max_{x \in [-B,B]} |p_d(x) - \sigma(x)| < \epsilon
    \]
    For $d = 7$, $B = 8$, $\epsilon < 10^{-3}$.
\end{definition}

\subsection{Persian Text Encoding}

Persian text presents specific challenges for FHE:

\begin{enumerate}
    \item \textbf{Right-to-left script.}  Token embeddings must preserve
    RTL ordering within CKKS slots.
    \item \textbf{Large vocabulary.}  The SentencePiece vocabulary for
    Persian contains 32K tokens; embedding lookup is implemented as
    encrypted matrix-vector multiplication.
    \item \textbf{Diacritics and Ezafe.}  Subword tokenisation must
    handle optional diacritical marks without information leakage
    through token count.
\end{enumerate}

\subsection{Inference Pipeline}

\begin{enumerate}
    \item User encodes input tokens as one-hot vectors in CKKS slots.
    \item User encrypts the packed vector under their CKKS public key.
    \item The inference node applies the model weights (in plaintext) to
    the encrypted input via homomorphic matrix multiplication.
    \item Between layers, the node performs rescaling and applies
    polynomial activations.
    \item The final logits (encrypted) are returned to the user for
    local decryption and argmax.
\end{enumerate}

\begin{theorem}[Inference Privacy]
    The inference node learns nothing about the input text, intermediate
    activations, or output logits beyond what is revealed by the
    ciphertext size and number of evaluation rounds.
\end{theorem}

\subsection{Performance}

\begin{table}[h]
\centering
\caption{Private inference benchmarks (4-layer transformer, 128 tokens).}
\label{tab:nlp_bench}
\begin{tabular}{@{}lrr@{}}
\toprule
\textbf{Configuration} & \textbf{Latency (s)} & \textbf{Accuracy Drop} \\
\midrule
Plaintext baseline & 0.02 & --- \\
CKKS, CPU          & 142.3 & 1.2\% \\
CKKS, GPU          & 18.7  & 1.2\% \\
\bottomrule
\end{tabular}
\end{table}

The 1.2\% accuracy drop results from polynomial activation approximation
and CKKS precision limits.  For classification tasks (sentiment, topic),
this is acceptable.

% ============================================================================
\section{Hybrid Pipeline and Scheme Switching}
\label{sec:hybrid}

Some applications require both exact (TFHE) and approximate (CKKS)
computation.  For example, a governance proposal may require exact
vote tallying (TFHE) combined with approximate sentiment analysis of
discussion text (CKKS).

\subsection{Scheme Switching}

\begin{definition}[TFHE-to-CKKS Switch]
    Given a TFHE ciphertext $c_{\mathrm{TFHE}}$ encrypting bit $\mu$,
    the scheme switch produces a CKKS ciphertext
    $c_{\mathrm{CKKS}}$ encrypting $\mu \in \{0.0, 1.0\}$ under
    a compatible CKKS key, via:
    \begin{enumerate}
        \item Functional bootstrap $c_{\mathrm{TFHE}}$ to extract
        $\mu$ into an RLWE ciphertext.
        \item Key-switch from the bootstrapping key space to the target
        CKKS key.
        \item Rescale to the CKKS encoding precision.
    \end{enumerate}
\end{definition}

The reverse switch (CKKS-to-TFHE) applies sign extraction followed by
modulus switching.

\subsection{Pipeline Composition}

\begin{enumerate}
    \item Encrypted search (TFHE) identifies matching messages.
    \item Scheme switch converts match indicators to CKKS.
    \item CKKS inference runs sentiment classification on matched messages.
    \item Results returned encrypted; user decrypts locally.
\end{enumerate}

% ============================================================================
\section{GPU Acceleration}
\label{sec:gpu}

FHE bootstrapping dominates computation cost.  We target GPU acceleration
on two fronts:

\subsection{TFHE Bootstrapping on GPU}

The TFHE bootstrap operation consists of:
\begin{enumerate}
    \item Blind rotation: $N$ external products, each an NTT-based
    polynomial multiplication.
    \item Key switching: matrix-vector multiplication in $\mathbb{Z}_q$.
\end{enumerate}

Both operations parallelise across GPU threads.  We adapt the cuFHE
approach~\cite{dai2018cufhe} with the following modifications for Pars:

\begin{itemize}
    \item Batch multiple independent gates into a single kernel launch.
    \item Use shared memory for the bootstrapping key (fits in 48\,KB
    for $n = 630$).
    \item Pipeline NTT with butterfly stages across warps.
\end{itemize}

\begin{theorem}[GPU Speedup]
    For $m$ independent TFHE gates, GPU bootstrapping achieves speedup
    $S = \min(m, C) \cdot \alpha$ over single-core CPU, where $C$ is
    the GPU core count and $\alpha \approx 3\text{--}5\times$ is the
    per-gate NTT speedup from dedicated hardware.
\end{theorem}

\subsection{CKKS on GPU}

CKKS homomorphic multiplication requires two NTTs on degree-$N$
polynomials.  For $N = 16384$, each NTT processes $2^{14}$ elements---a
natural fit for GPU parallelism.

The Pars GPU pipeline (PNP-011) provides the NTT infrastructure shared
with post-quantum cryptography (PNP-004).

% ============================================================================
\section{Security Analysis}
\label{sec:security}

\subsection{Cryptographic Assumptions}

\begin{assumption}[Ring-LWE]
    For dimension $N$, modulus $q$, and Gaussian error distribution
    $\chi$, the Ring-LWE problem is computationally hard: no
    polynomial-time algorithm distinguishes $(a, a \cdot s + e)$ from
    uniform over $R_q \times R_q$, where $R_q = \mathbb{Z}_q[X]/(X^N+1)$.
\end{assumption}

Both TFHE and CKKS security reduce to Ring-LWE.  Our parameters follow
the Homomorphic Encryption Standard~\cite{hes2018}, targeting $\lambda
\geq 128$ bits.

\subsection{Side-Channel Resistance}

FHE evaluation on untrusted nodes must resist timing and power
side-channels:

\begin{itemize}
    \item \textbf{Constant-time NTT.}  All polynomial operations use
    constant-time butterfly networks.
    \item \textbf{Oblivious memory access.}  Bootstrapping key access
    patterns are data-independent.
    \item \textbf{Ciphertext padding.}  All ciphertexts are padded to
    uniform size to prevent length-based information leakage.
\end{itemize}

\subsection{Post-Quantum Security}

Both TFHE and CKKS are based on lattice problems and are believed to
resist quantum attacks.  The parameter sets in Table~\ref{tab:params}
provide 128-bit post-quantum security under the methodology of
PNP-004~\cite{pnp004_pq}.

% ============================================================================
\section{Integration with Pars Infrastructure}
\label{sec:integration}

\subsection{Session Layer (PNP-003)}

The encrypted search index is stored alongside E2E-encrypted messages
on Session daemons.  The TFHE public key is distributed via the same
prekey bundle mechanism used for session establishment.

\subsection{Governance (PNP-005)}

FHE voting replaces the current mix-net approach for ballot tallying.
The threshold TFHE key is managed by the governance committee using
the threshold signing protocol (PNP-010).

\subsection{Distributed Inference}

The Pars distributed inference framework currently exposes plaintext
queries to inference nodes.  CKKS inference wraps the existing model
serving pipeline: the model weights remain in plaintext on the node,
but all user data is encrypted.

\subsection{Lux Network}

Pars inherits FHE research from the Lux ecosystem (lux-tfhe), sharing
parameter sets and GPU kernels.  Lux's FHE-enabled confidential smart
contracts provide the L1 settlement layer for Pars FHE computations.

% ============================================================================
\section{Performance Summary}
\label{sec:performance}

\begin{table}[h]
\centering
\caption{End-to-end latency for FHE applications on Pars.}
\label{tab:e2e}
\begin{tabular}{@{}lrrr@{}}
\toprule
\textbf{Application} & \textbf{CPU (s)} & \textbf{GPU (s)} & \textbf{Overhead} \\
\midrule
Search (1K msgs)      & 8.2   & 1.4   & $100\times$ \\
Vote tally (10K)      & 24.1  & 3.8   & $190\times$ \\
NLP inference (128 tok) & 142.3 & 18.7  & $935\times$ \\
\bottomrule
\end{tabular}
\end{table}

The overhead relative to plaintext computation is substantial but
acceptable for the security properties gained.  Encrypted search at
1.4\,s GPU latency is within interactive thresholds.  Vote tallying is a
batch operation with no real-time constraint.  NLP inference at 18.7\,s
is viable for asynchronous queries (translation, summarisation) though
not for real-time conversation.

% ============================================================================
\section{Related Work}
\label{sec:related}

\paragraph{Concrete.}  Zama's Concrete framework~\cite{zama2022concrete}
provides a TFHE compiler for Rust and Python.  Pars builds on the same
TFHE parameter sets but adds diaspora-specific optimisations (Persian
text handling, mesh-network key distribution).

\paragraph{SEAL.}  Microsoft SEAL~\cite{seal2022} implements BFV and CKKS.
Pars uses CKKS parameters compatible with SEAL but deploys on
decentralised inference nodes rather than centralised cloud.

\paragraph{HELib.}  IBM's HELib~\cite{halevi2014helib} pioneered
practical BGV.  Pars requires TFHE's gate-level bootstrapping for
exact circuits (search, voting), which HELib does not support.

\paragraph{Lux-TFHE.}  The Lux network's FHE infrastructure provides
shared NTT kernels and parameter sets.  Pars extends these with
application-layer protocols specific to diaspora use cases.

% ============================================================================
\section{Conclusion}
\label{sec:conclusion}

We have presented a comprehensive FHE framework for the Pars Network
that enables encrypted search over messaging archives, confidential
governance vote tallying, and private NLP inference on Persian text.
By combining TFHE for exact circuits with CKKS for approximate
arithmetic, and accelerating both via GPU, we achieve practical
latencies for all three applications.

The integration with Pars's existing session layer, governance
protocols, and distributed inference framework provides end-to-end
encrypted computation pipelines.  No plaintext user data is exposed
to any untrusted node at any point in the computation.

Future work includes: (1)~multi-key FHE to eliminate the threshold
key setup, (2)~FHE-friendly model architectures that reduce
multiplicative depth, and (3)~client-side hardware acceleration
for mobile devices commonly used by diaspora communities.

% ============================================================================
\begin{thebibliography}{99}

\bibitem{gentry2009fhe}
C.~Gentry,
``Fully Homomorphic Encryption Using Ideal Lattices,''
\emph{STOC}, pp.~169--178, 2009.

\bibitem{chillotti2020tfhe}
I.~Chillotti, N.~Gama, M.~Georgieva, and M.~Izabach\`{e}ne,
``TFHE: Fast Fully Homomorphic Encryption over the Torus,''
\emph{Journal of Cryptology}, vol.~33, pp.~34--91, 2020.

\bibitem{cheon2017ckks}
J.~H.~Cheon, A.~Kim, M.~Kim, and Y.~Song,
``Homomorphic Encryption for Arithmetic of Approximate Numbers,''
\emph{ASIACRYPT}, pp.~409--437, 2017.

\bibitem{hes2018}
M.~Albrecht et al.,
``Homomorphic Encryption Standard,''
\emph{HomomorphicEncryption.org}, 2018.

\bibitem{dai2018cufhe}
W.~Dai and B.~Sunar,
``cuFHE: A GPU Library for Fully Homomorphic Encryption,''
\emph{IACR ePrint 2017/007}, 2018.

\bibitem{zama2022concrete}
Zama,
``Concrete: TFHE Compiler for Secure Computation,''
2022.

\bibitem{seal2022}
Microsoft Research,
``Microsoft SEAL (release 4.0),''
2022.

\bibitem{halevi2014helib}
S.~Halevi and V.~Shoup,
``Algorithms in HElib,''
\emph{CRYPTO}, pp.~554--571, 2014.

\bibitem{pnp003}
Cyrus Pahlavi, Pars Team,
``PNP-003: Pars Session Protocol,''
2026.

\bibitem{pnp004_pq}
Cyrus Pahlavi, Pars Team,
``PNP-004: Post-Quantum Cryptographic Standards for Pars Network,''
2026.

\bibitem{pnp005}
Cyrus Pahlavi, Pars Team,
``PNP-005: Coercion-Resistant Protocols for Governance Under
Authoritarian Surveillance,''
2026.

\bibitem{pars_distributed_inference}
Cyrus Pahlavi, Pars Team,
``Distributed Inference for Persian Language Models on Pars Network,''
2026.

\end{thebibliography}

\end{document}
